% Answer Key: Prompt Engineering (Weeks 6 & 8)
\documentclass[11pt,a4paper]{article}
\usepackage[utf8]{inputenc}
\usepackage{geometry}
\usepackage{enumitem}
\geometry{margin=1in}
\begin{document}
\begin{center}
  {\LARGE Answer Key \textemdash{} Prompt Engineering Quiz 3}\\[6pt]
\end{center}
\vspace{0.5cm}
\begin{enumerate}[leftmargin=*, label=\textbf{\#\arabic*}]

  \item \textbf{B} \textemdash{} Constraint Violation. Ignoring negative constraints like "no lists" is a classic constraint failure.
  \item \textbf{B} \textemdash{} To ask the model to evaluate its own output against a rubric before revising. This separates the "creation" step from the "editing" step.
  \item \textbf{C} \textemdash{} Fragility and inability to fix edge cases. If you don't understand the code generated by "vibes," you cannot maintain it.
  \item \textbf{B} \textemdash{} Using a strong model (e.g., GPT-4) to grade the outputs of other models. This scales evaluation without human fatigue.
  \item \textbf{B} \textemdash{} Evaluate. This phase checks outputs against requirements before refinement or deployment.
  \item \textbf{B} \textemdash{} It is significantly easier to generate text than to verify it. This asymmetry creates the "Evaluation Gap."
  \item \textbf{D} \textemdash{} Human-like. The 3 H's are Helpful, Honest, and Harmless.
  \item \textbf{B} \textemdash{} A curated set of inputs and "perfect" outputs used as a benchmark for testing.
  \item \textbf{C} \textemdash{} Syntax or JSON format compliance. Hard metrics are deterministic and binary (pass/fail).
  \item \textbf{B} \textemdash{} Rate longer answers as better, even if they contain fluff. This is a common bias in both humans and LLMs.
  \item \textbf{B} \textemdash{} Risk-based sampling. This focuses evaluation effort on high-stakes or new areas.
  \item \textbf{C} \textemdash{} Decide which of two responses (A or B) is better. This is often more reliable than absolute scoring.

\end{enumerate}
\end{document}
